% ============================================================
%  Cahier des charges -- PFA-21
%  Compatible pdfLaTeX, XeLaTeX et LuaLaTeX (2 passes)
% ============================================================
\documentclass[11pt,a4paper]{article}

\usepackage{iftex}
\ifPDFTeX
  \usepackage[utf8]{inputenc}
  \usepackage[T1]{fontenc}
  \usepackage{lmodern}
  \usepackage[french]{babel}
  \newcommand{\code}[1]{{\NoAutoSpacing\texttt{\small #1}}}
\else
  \usepackage{fontspec}
  \setmainfont{Latin Modern Roman}
  \setsansfont{Latin Modern Sans}
  \setmonofont{Latin Modern Mono}
  \usepackage{polyglossia}
  \setmainlanguage{french}
  \newcommand{\code}[1]{{\texttt{\small #1}}}
\fi
\usepackage[a4paper,margin=2.5cm,headheight=15pt]{geometry}
\usepackage{microtype}
\usepackage{xcolor}
\usepackage{graphicx}
\usepackage{array}
\usepackage{booktabs}
\usepackage{longtable}
\usepackage{enumitem}
\usepackage{titlesec}
\usepackage{fancyhdr}
\usepackage{listings}
\usepackage{tcolorbox}
\usepackage{setspace}
\usepackage[hidelinks]{hyperref}

% ---------- Couleurs ----------
\definecolor{bleu}{HTML}{1F3A5F}
\definecolor{bleuclair}{HTML}{E8EEF6}
\definecolor{gris}{HTML}{F4F4F4}
\definecolor{vert}{HTML}{2E7D32}
\definecolor{rouge}{HTML}{C62828}

% ---------- Raccourcis ----------
\newcommand{\PASS}{\textbf{\textcolor{vert}{PASS}}}
\newcommand{\FAIL}{\textbf{\textcolor{rouge}{FAIL}}}
\newcommand{\WARN}{\textbf{\textcolor{orange}{WARN}}}

% ---------- Titres ----------
\titleformat{\section}{\Large\bfseries\color{bleu}}{\thesection.}{0.6em}{}
\titleformat{\subsection}{\large\bfseries\color{bleu}}{\thesubsection}{0.6em}{}
\titleformat{\subsubsection}{\normalsize\bfseries\color{bleu}}{\thesubsubsection}{0.6em}{}
\titlespacing*{\section}{0pt}{1.6em}{0.9em}
\titlespacing*{\subsection}{0pt}{1.2em}{0.5em}

% ---------- En-tête / pied de page ----------
\pagestyle{fancy}
\fancyhf{}
\fancyhead[L]{\small\color{bleu}PFA-21 --- Audit statique d'une application Android}
\fancyhead[R]{\small\color{bleu}Cahier des charges}
\fancyfoot[L]{\small EMSI --- Cybersécurité et Réseaux}
\fancyfoot[R]{\small Page \thepage}
\renewcommand{\headrulewidth}{0.4pt}
\renewcommand{\footrulewidth}{0.4pt}
\fancypagestyle{plain}{\fancyhf{}\renewcommand{\headrulewidth}{0pt}\renewcommand{\footrulewidth}{0pt}}

% ---------- Listes et tableaux ----------
\setlist[itemize]{leftmargin=1.4em,itemsep=2pt,topsep=4pt}
\setlist[enumerate]{leftmargin=1.6em,itemsep=2pt,topsep=4pt}
\renewcommand{\arraystretch}{1.3}
\setlength{\parindent}{0pt}
\setlength{\parskip}{6pt}
\setlength{\emergencystretch}{2em}
\newcommand{\hd}[1]{\textcolor{white}{\textbf{#1}}}
\usepackage{colortbl}
\newcolumntype{L}[1]{>{\raggedright\arraybackslash}p{#1}}

% ---------- Code ----------
\lstset{
  basicstyle=\ttfamily\footnotesize,
  backgroundcolor=\color{gris},
  frame=single,
  rulecolor=\color{bleu!40},
  breaklines=true,
  columns=fullflexible,
  keepspaces=true,
  xleftmargin=0.5em,xrightmargin=0.5em,
  aboveskip=8pt,belowskip=8pt
}

% ---------- Encadré ----------
\newtcolorbox{remarque}[1][Remarque]{
  colback=bleuclair,colframe=bleu,fonttitle=\bfseries,title=#1,
  boxrule=0.5pt,arc=2pt,left=6pt,right=6pt,top=4pt,bottom=4pt
}

\hypersetup{
  pdftitle={Cahier des charges PFA-21},
  pdfauthor={BENLYAMANI Ziad, FERDOUS Larbi Sifeddine},
  pdfsubject={Audit statique très ciblé d'une application Android}
}

\begin{document}

% ============================================================
%  PAGE DE GARDE
% ============================================================
\begin{titlepage}
\thispagestyle{empty}
\centering
\setstretch{1.15}

% Logo : décommenter et placer le fichier à côté du .tex
% \includegraphics[height=2.4cm]{logo_emsi.png}\par

{\large\color{bleu}\textbf{École Marocaine des Sciences de l'Ingénieur (EMSI)}}\par
{\normalsize Ingénierie Informatique et Réseaux --- Option Cybersécurité et Réseaux}\par

\vspace{2.2cm}
{\color{bleu}\rule{\textwidth}{1.6pt}}\par
\vspace{0.6cm}
{\Large\bfseries\color{bleu} PROJET PFA}\par
\vspace{0.4cm}
{\huge\bfseries Audit statique très ciblé\\[4pt] d'une application Android}\par
\vspace{0.6cm}
{\color{bleu}\rule{\textwidth}{1.6pt}}\par

\vspace{1.2cm}
{\Large\bfseries Cahier des charges}\par
\vspace{0.3cm}
{\large Projet PFA-21}\par

\vfill

\begin{tabular}{p{7.2cm}p{7.2cm}}
\textbf{\color{bleu}\large Réalisé par :} & \textbf{\color{bleu}\large Encadré par :}\\[8pt]
BENLYAMANI Ziad & Pr. OUAGUID Abdellah\\
FERDOUS Larbi Sifeddine &\\
\end{tabular}

\vfill

{\normalsize Année universitaire 2026--2027}\par
{\normalsize Version 1.0 --- Novembre 2026}\par
\end{titlepage}

% ============================================================
%  SOMMAIRE
% ============================================================
\pagenumbering{roman}
\tableofcontents
\clearpage
\pagenumbering{arabic}

% ============================================================
\section{Présentation du projet}
% ============================================================

\subsection{Contexte}
Une équipe fictive développe une application Android interne permettant de consulter le programme d'un événement (« EventSchedule »). Avant toute évolution, elle souhaite s'assurer que le fichier \code{AndroidManifest.xml} ne contient aucun choix risqué : permission excessive, composant exposé ou configuration de sécurité imprudente.

\subsection{Problématique}
Comment vérifier de façon automatique, reproductible et justifiée que le manifeste d'une application Android respecte un socle minimal de sécurité, puis démontrer qu'une correction modifie bien le constat ?

\subsection{Objectif général}
Détecter puis corriger un choix de configuration Android non sécurisé, par un audit statique outillé en Python, avec preuves avant/après.

\subsection{Objectifs spécifiques}
\begin{itemize}
  \item \textbf{OS1} : sélectionner un projet Android libre ou pédagogique autorisé (ou en créer un minimal).
  \item \textbf{OS2} : définir une grille de 4 contrôles de manifeste, chacun avec règle, sévérité, justification et référence.
  \item \textbf{OS3} : développer un script Python lisant le manifeste et produisant un état (\PASS{} / \WARN{} / \FAIL{}) et une justification par contrôle.
  \item \textbf{OS4} : introduire volontairement un paramètre trop permissif, le détecter, le corriger, relancer l'audit.
  \item \textbf{OS5} : automatiser \textit{setup / start / test / stop} et documenter les actions non automatisables.
  \item \textbf{OS6} : permettre à chaque membre de redéployer et valider le MVP sur sa propre machine.
\end{itemize}

\subsection{Indicateurs de réussite}
\begin{center}
\begin{tabular}{L{9.2cm}L{5.4cm}}
\rowcolor{bleu}\hd{Indicateur} & \hd{Cible}\\
Contrôles implémentés avec état et justification & 4 / 4\\
Constat \FAIL{} sur le manifeste « avant » & au moins 1 (paramètre ciblé)\\
Constat \PASS{} sur le même contrôle après correction & 100\,\%\\
Tests de recette réussis avec preuve & 3 / 3\\
Redéploiement réussi sur le poste de chaque membre & 100\,\%\\
Tests automatisés Python au vert & 100\,\%\\
\end{tabular}
\end{center}

% ============================================================
\section{Périmètre}
% ============================================================

\subsection{Inclus (MVP)}
\begin{itemize}
  \item Analyse statique du seul \code{AndroidManifest.xml} (et des fichiers XML qu'il référence directement, par exemple \code{network\_security\_config.xml}).
  \item 4 contrôles de manifeste et de permissions.
  \item Un manifeste « avant » (volontairement vulnérable) et un manifeste « après » (corrigé).
  \item Script Python, tests automatisés, rapports (console et JSON), preuves.
\end{itemize}

\subsection{Exclus}
\begin{itemize}
  \item Publication sur un magasin d'applications, usage d'un téléphone personnel ou réel.
  \item Pentest mobile avancé, HarmonyOS natif, certification MASVS.
  \item Analyse du code Java/Kotlin, des bibliothèques tierces et du trafic réseau.
\end{itemize}

\subsection{Extensions (après validation du MVP, avec accord de l'encadrant)}
\begin{itemize}
  \item Analyse complémentaire avec MobSF (conteneur Docker) et comparaison avec le script.
  \item Compilation de l'APK (Gradle) pour vérifier le manifeste fusionné.
  \item Test dynamique sur émulateur.
\end{itemize}

\subsection{Hypothèses et contraintes}
\begin{itemize}
  \item Tout le travail est local (conteneur ou machine personnelle), sans appareil réel.
  \item Seules des applications libres, pédagogiques ou créées par l'équipe sont analysées.
  \item Les exigences peuvent évoluer sur décision de l'encadrant.
  \item Chaque membre maîtrise l'intégralité du projet.
\end{itemize}

% ============================================================
\section{Parties prenantes et rôles}
% ============================================================
\begin{center}
\begin{tabular}{L{3.6cm}L{11cm}}
\rowcolor{bleu}\hd{Rôle} & \hd{Responsabilités}\\
Encadrant & Valide le sujet, le cahier des charges, les extensions et la recette\\
Membre A & Vérifications, références (OWASP, Android), définition des règles et de la grille\\
Membre B & Correction du paramètre, automatisation (scripts), collecte des preuves\\
Membre C (si trinôme) & Tests, scripts de validation, documentation\\
\end{tabular}
\end{center}
En binôme, les tâches du membre C sont réparties entre A et B.

% ============================================================
\section{Environnement et outillage}
% ============================================================
\begin{center}
\begin{tabular}{L{3.6cm}L{11cm}}
\rowcolor{bleu}\hd{Élément} & \hd{Choix}\\
Système & Linux, Windows ou macOS avec Docker, ou machine virtuelle Linux\\
Conteneur & Image Python 3.11+ (Dockerfile fourni)\\
Langage & Python 3 (bibliothèque standard \code{xml.etree.ElementTree}, option \code{defusedxml})\\
Tests & \code{pytest}\\
Versionnage & Git\\
Optionnel & Android Studio (consultation/édition), MobSF (extension)\\
\end{tabular}
\end{center}

% ============================================================
\section{Application cible}
% ============================================================

\subsection{Critères de sélection}
\begin{itemize}
  \item Licence libre ou usage pédagogique explicite ; code source disponible ; manifeste lisible.
  \item Petite taille, sans dépendance à un appareil.
  \item Possibilité de modifier le manifeste sans casser le projet.
\end{itemize}

\subsection{Choix recommandé}
Créer (ou adapter) un mini-projet « EventSchedule » : une activité principale, un service ou récepteur interne, un fournisseur de contenu et quelques permissions. Avantage : contrôle total du manifeste et aucun problème de licence. Alternative : une application volontairement vulnérable open source (par exemple DIVA), après vérification de sa licence et accord de l'encadrant.

\subsection{Paramètre volontairement permissif (cas d'étude)}
Un composant interne (par exemple un \code{Service} ou un \code{BroadcastReceiver} de synchronisation du programme) déclaré avec \code{android:exported="true"} sans permission de protection. Les trois autres contrôles restent conformes dans le manifeste « avant », afin que le constat évolue uniquement sur le point ciblé.

% ============================================================
\section{Grille des 4 contrôles}
% ============================================================
\begin{longtable}{@{}L{0.8cm}L{2.4cm}L{2.9cm}L{4.6cm}L{2.7cm}@{}}
\rowcolor{bleu}\color{white}\bfseries ID & \color{white}\bfseries Contrôle & \color{white}\bfseries Élément analysé & \color{white}\bfseries Règle de décision & \color{white}\bfseries Références\\
\endhead
C1 & Permissions excessives & \code{<uses-permission>} & \FAIL{} si une permission dangereuse non justifiée figure hors des besoins déclarés ; \WARN{} si sensible mais plausible & Android (permissions dangereuses), CWE-250, MASVS-PLATFORM\\
C2 & Composants exportés non protégés & \code{<activity>}, \code{<service>}, \code{<receiver>}, \code{<provider>} & \FAIL{} si \code{exported="true"} (explicite, ou implicite via \code{intent-filter} avec \code{targetSdk} $<$ 31) sans \code{android:permission}, hors activité LAUNCHER & CWE-926, MASVS-PLATFORM\\
C3 & Drapeaux applicatifs & \code{android:debuggable}, \code{android:allowBackup} & \FAIL{} si \code{debuggable="true"} ; \WARN{} si \code{allowBackup="true"} (ou absent selon la cible) & CWE-489, MASVS-STORAGE\\
C4 & Trafic en clair & \code{usesCleartextTraffic}, \code{networkSecurityConfig} & \FAIL{} si \code{usesCleartextTraffic="true"} ou configuration autorisant le clair globalement ; \PASS{} sinon & CWE-319, MASVS-NETWORK\\
\end{longtable}

\begin{remarque}[À valider]
Le choix des 4 contrôles est adaptable avec l'accord de l'encadrant. Les références (CWE, MASVS) sont indicatives et doivent être vérifiées et citées avec leur source dans le dossier final.
\end{remarque}

Pour chaque contrôle, la grille documente : objectif, risque, condition \FAIL{}/\WARN{}/\PASS{}, justification type, remédiation recommandée et faux positifs connus.

% ============================================================
\section{Exigences fonctionnelles}
% ============================================================
\begin{longtable}{@{}L{1.1cm}L{5.9cm}L{2.1cm}L{4.5cm}@{}}
\rowcolor{bleu}\color{white}\bfseries ID & \color{white}\bfseries Exigence & \color{white}\bfseries Priorité & \color{white}\bfseries Critère d'acceptation\\
\endhead
EF1 & Sélectionner et documenter l'application cible (source, licence, version) & Obligatoire & Fiche de cible validée par l'encadrant\\
EF2 & Signaler le paramètre volontairement permissif & Obligatoire & Contrôle concerné en \FAIL{} sur le manifeste « avant »\\
EF3 & Le constat évolue après modification du manifeste & Obligatoire & Même contrôle en \PASS{} sur le manifeste « après »\\
EF4 & Chaque contrôle fournit un état et une justification & Obligatoire & 4 lignes (ID, état, justification) par rapport\\
EF5 & Le script lit le manifeste via un chemin passé en argument & Obligatoire & \code{python audit.py chemin/AndroidManifest.xml} fonctionne\\
EF6 & Sortie console lisible et sortie JSON & Obligatoire & Fichier JSON valide, produit à la demande\\
EF7 & Code de sortie exploitable (0 conforme, 1 constat \FAIL{}, 2 erreur) & Obligatoire & Vérifié par test automatisé\\
EF8 & Comparaison avant/après & Obligatoire & Commande ou rapport listant les contrôles qui changent\\
EF9 & Gestion des erreurs (fichier absent, XML invalide) & Obligatoire & Message clair, code 2, pas de trace brute\\
EF10 & Extension MobSF & Optionnel & Rapport MobSF comparé à celui du script\\
\end{longtable}

% ============================================================
\section{Exigences non fonctionnelles}
% ============================================================
\begin{itemize}
  \item \textbf{Reproductibilité} : déploiement identique sur chaque poste via conteneur et scripts versionnés.
  \item \textbf{Sécurité du traitement} : parseur XML durci (pas d'entités externes), aucune exécution de code issu de la cible, aucun accès réseau requis à l'exécution.
  \item \textbf{Lisibilité} : code commenté, PEP~8, une fonction par contrôle.
  \item \textbf{Extensibilité} : ajout d'un contrôle sans modifier le noyau (registre de contrôles).
  \item \textbf{Traçabilité} : chaque constat renvoie à l'élément XML et à la règle appliquée.
  \item \textbf{Éthique et légalité} : cibles autorisées uniquement.
  \item \textbf{Documentation} : toute commande est documentée et testée à froid.
\end{itemize}

% ============================================================
\section{Spécification du script d'audit}
% ============================================================

\subsection{Interface}
\begin{lstlisting}
python audit.py <manifest.xml> [--json rapport.json] [--compare autre.xml]
\end{lstlisting}

\subsection{Fonctionnement}
\begin{enumerate}
  \item Charger et valider le XML (gestion de l'espace de noms \code{android}).
  \item Extraire : \code{targetSdkVersion} (si présent), permissions, composants et leurs attributs, attributs de \code{<application>}.
  \item Exécuter les contrôles C1 à C4 (fonctions indépendantes retournant état, justification, élément concerné).
  \item Agréger, afficher, exporter en JSON et retourner le code de sortie.
\end{enumerate}

\subsection{Format de sortie JSON (exemple)}
\begin{lstlisting}
{
  "fichier": "manifests/avant/AndroidManifest.xml",
  "controles": [
    {"id": "C2", "etat": "FAIL", "element": "service .SyncService",
     "justification": "exported=true sans permission"}
  ],
  "synthese": {"PASS": 3, "WARN": 0, "FAIL": 1}
}
\end{lstlisting}

\subsection{Points techniques à traiter}
\begin{itemize}
  \item Valeur par défaut de \code{exported} selon la version cible (comportement différent avant et après Android~12).
  \item Attribut absent versus explicitement à \code{false}.
  \item Activité LAUNCHER exportée légitimement : éviter le faux positif.
  \item Manifeste fusionné (APK) versus manifeste source : le préciser dans le rapport.
\end{itemize}

% ============================================================
\section{Automatisation : setup / start / test / stop}
% ============================================================
\begin{center}
\begin{tabular}{L{2.8cm}L{3.6cm}L{8.2cm}}
\rowcolor{bleu}\hd{Commande} & \hd{Rôle} & \hd{Détail}\\
\code{make setup} & Prépare l'environnement & Construit l'image Docker, installe \code{pytest}, vérifie les versions\\
\code{make start} & Lance l'audit & Exécute le script sur le manifeste « avant » et génère les rapports\\
\code{make test} & Valide & Lance \code{pytest} puis le scénario avant/après et vérifie le changement de constat\\
\code{make stop} & Nettoie & Arrête et supprime le conteneur, purge les fichiers temporaires\\
\end{tabular}
\end{center}
Des équivalents shell/PowerShell sont fournis si \code{make} est indisponible.

\textbf{Actions non automatisables (à documenter) :} choix et validation de l'application par l'encadrant, installation de Docker ou d'Android Studio, captures d'écran des preuves visuelles, interprétation métier des constats, rédaction de la justification des \WARN{}.

% ============================================================
\section{Recette obligatoire}
% ============================================================
\begin{longtable}{@{}L{0.9cm}L{3.3cm}L{3.3cm}L{2.2cm}L{3.6cm}@{}}
\rowcolor{bleu}\color{white}\bfseries Test & \color{white}\bfseries Action & \color{white}\bfseries Attendu & \color{white}\bfseries Résultat & \color{white}\bfseries Preuve\\
\endhead
T1 & Exécuter les 4 contrôles sur le manifeste « avant » & 4 lignes (état et justification) & à compléter & Capture console et JSON\\
T2 & Vérifier la détection du paramètre permissif & Contrôle ciblé en \FAIL{}, élément nommé & à compléter & Rapport « avant »\\
T3 & Corriger le paramètre puis relancer & Contrôle ciblé en \PASS{}, autres inchangés & à compléter & Rapport « après », diff Git, journal de \code{make test}\\
\end{longtable}

Chaque test est consigné dans une fiche : identifiant, date, exécutant, machine, action, attendu, résultat, preuve (chemin du fichier), verdict.

% ============================================================
\section{Livrables}
% ============================================================
\begin{enumerate}
  \item Cahier des charges (le présent document, version finale).
  \item Grille des 4 contrôles complète.
  \item Projet ou manifeste cible : versions « avant » et « après ».
  \item Script \code{audit.py}, tests \code{pytest}, scripts \textit{setup/start/test/stop}, Dockerfile, Makefile.
  \item Dossier de preuves (captures, rapports JSON, journaux) pour chaque membre.
  \item Rapport de recette et rapport final.
  \item Support de soutenance et démonstration reproductible.
\end{enumerate}

\subsection{Structure du dépôt}
\begin{lstlisting}
pfa21/
|-- app/manifests/{avant,apres}/AndroidManifest.xml
|-- src/audit.py
|-- tests/
|-- scripts/{setup,start,test,stop}
|-- docs/{cahier_des_charges,grille,recette,rapport}
|-- preuves/
|-- Dockerfile
|-- Makefile
`-- README.md
\end{lstlisting}

% ============================================================
\section{Risques et parades}
% ============================================================
\begin{center}
\begin{tabular}{L{4.4cm}L{2.2cm}L{7.8cm}}
\rowcolor{bleu}\hd{Risque} & \hd{Impact} & \hd{Parade}\\
Application cible inadaptée ou licence floue & Retard & Projet maison « EventSchedule » en secours\\
Faux positifs (LAUNCHER, valeurs par défaut) & Mauvais constat & Règles documentées, tests de cas limites\\
Environnement non reproductible sur un poste & Recette non validée & Docker, guide de dépannage, test à froid au sprint 5\\
Dérive vers MobSF/APK avant le MVP & Retard & Extensions gelées jusqu'à validation du MVP\\
Membre ne maîtrisant pas tout & Pénalité en soutenance & Revue croisée et démonstration individuelle\\
\end{tabular}
\end{center}

% ============================================================
\section{Critères d'acceptation globale}
% ============================================================
\begin{itemize}
  \item Les 3 exigences du MVP et les 3 tests de recette sont validés avec preuves.
  \item Un tiers reproduit \textit{setup $\rightarrow$ start $\rightarrow$ test $\rightarrow$ stop} sans aide.
  \item Le dossier est cohérent : chaque constat est relié à une règle et à une référence.
  \item Aucun secret, donnée personnelle ou appareil réel n'est impliqué.
\end{itemize}

% ============================================================
\section{Planification en six sprints}
% ============================================================
Durée indicative : une semaine par sprint (six semaines), ajustable selon le calendrier de l'encadrant. Chaque sprint se termine par une revue et une démonstration.

\subsection{Sprint 1 --- Cadrage et mise en place du laboratoire}
\textbf{Objectif :} valider le sujet et disposer d'un environnement fonctionnel sur chaque poste.
\begin{itemize}
  \item Valider le sujet, le périmètre et ce cahier des charges avec l'encadrant ; répartir les rôles.
  \item Choisir l'application cible, vérifier sa licence, rédiger la fiche de cible (source, version, SDK cible).
  \item Créer le dépôt Git, la structure de dossiers et les conventions (branches, commits).
  \item Écrire le Dockerfile et un premier \code{make setup}.
  \item Récupérer ou créer le manifeste de base et le versionner.
  \item Établir le backlog, la matrice des risques et le planning.
\end{itemize}
\textbf{Livrables :} fiche de cible, dépôt initialisé, conteneur opérationnel, backlog.\\
\textbf{Critères d'acceptation :} \code{make setup} réussit sur chaque poste ; l'encadrant valide la cible.

\subsection{Sprint 2 --- Référentiel et grille des 4 contrôles}
\textbf{Objectif :} disposer d'une grille de contrôle précise, justifiée et validée.
\begin{itemize}
  \item Étudier la structure du manifeste et la documentation Android officielle (permissions, \code{exported}, sauvegarde, trafic en clair).
  \item Sélectionner les 4 contrôles et formaliser pour chacun : risque, règle \PASS{}/\WARN{}/\FAIL{}, sévérité, remédiation, faux positifs.
  \item Recenser les références (CWE, OWASP MASVS) avec sources vérifiées.
  \item Choisir le paramètre volontairement permissif et préparer les manifestes « avant » et « après ».
  \item Faire valider la grille par l'encadrant.
\end{itemize}
\textbf{Livrables :} grille v1, manifestes avant/après, tableau de références.\\
\textbf{Critères d'acceptation :} chaque contrôle a une règle sans ambiguïté ; la différence avant/après ne porte que sur le point ciblé.

\subsection{Sprint 3 --- Développement du script d'audit}
\textbf{Objectif :} produire un script fonctionnel qui applique les 4 contrôles.
\begin{itemize}
  \item Implémenter le chargement sécurisé du XML et la gestion des espaces de noms.
  \item Implémenter C1 à C4 sous forme de fonctions indépendantes avec registre.
  \item Gérer les cas limites : valeur par défaut de \code{exported}, LAUNCHER, attributs absents.
  \item Implémenter l'affichage console, l'export JSON et les codes de sortie.
  \item Gérer les erreurs (fichier absent, XML invalide, manifeste vide).
  \item Effectuer une relecture de code croisée.
\end{itemize}
\textbf{Livrables :} \code{audit.py} v1, exemples de rapports.\\
\textbf{Critères d'acceptation :} le script produit 4 états avec justification sur les deux manifestes ; le contrôle ciblé est \FAIL{} avant.

\subsection{Sprint 4 --- Tests automatisés et chaîne d'automatisation}
\textbf{Objectif :} fiabiliser le script et automatiser tout le cycle.
\begin{itemize}
  \item Écrire les tests \code{pytest} : un cas \PASS{} et un cas \FAIL{} par contrôle, cas limites, cas d'erreur.
  \item Écrire le test avant/après (le constat ciblé passe de \FAIL{} à \PASS{}, les autres restent identiques).
  \item Implémenter la comparaison \code{--compare}.
  \item Finaliser \code{make setup/start/test/stop} et leurs équivalents shell.
  \item Rédiger la liste des actions non automatisables.
  \item Mesurer la couverture et corriger les défauts.
\end{itemize}
\textbf{Livrables :} suite de tests, scripts de déploiement et de validation, guide d'exécution.\\
\textbf{Critères d'acceptation :} \code{make test} vert ; le cycle complet s'exécute sans intervention manuelle.

\subsection{Sprint 5 --- Correction, vérification et recette}
\textbf{Objectif :} exécuter officiellement la recette et collecter les preuves.
\begin{itemize}
  \item Appliquer la correction du paramètre sur le projet de test (commit dédié).
  \item Relancer l'audit et conserver les rapports avant/après et le diff Git.
  \item Exécuter T1, T2, T3 en renseignant action, attendu, résultat et preuve.
  \item Chaque membre redéploie et valide le MVP sur sa propre machine ; collecte de ses preuves.
  \item Traiter les anomalies relevées et rejouer les tests concernés.
\end{itemize}
\textbf{Livrables :} rapport de recette signé, dossier de preuves par membre.\\
\textbf{Critères d'acceptation :} 3 tests sur 3 validés sur chaque poste.

\subsection{Sprint 6 --- Documentation, extension et soutenance}
\textbf{Objectif :} livrer un dossier final cohérent et préparer la soutenance.
\begin{itemize}
  \item Rédiger le rapport final : contexte, méthode, grille, résultats, limites, leçons apprises.
  \item Mettre à jour le README (installation, usage, dépannage) et tester à froid sur une machine propre.
  \item Extension optionnelle, uniquement si le MVP est validé : MobSF en conteneur et comparaison avec le script.
  \item Préparer la démonstration (scénario avant $\rightarrow$ correction $\rightarrow$ après) et le support de soutenance.
  \item Répéter la soutenance avec questions croisées pour garantir la maîtrise individuelle.
  \item Geler la version (tag Git) et remettre les livrables.
\end{itemize}
\textbf{Livrables :} rapport final, dépôt versionné, support, démonstration.\\
\textbf{Critères d'acceptation :} reproduction par un tiers ; chaque membre explique et exécute l'ensemble du projet.

% ============================================================
\section{Synthèse du planning}
% ============================================================
\begin{center}
\begin{tabular}{L{1.6cm}L{6cm}L{6.6cm}}
\rowcolor{bleu}\hd{Sprint} & \hd{Thème} & \hd{Jalon de sortie}\\
S1 & Cadrage et laboratoire & Cible validée, environnement prêt\\
S2 & Référentiel et grille & Grille et manifestes validés\\
S3 & Script d'audit & 4 contrôles fonctionnels\\
S4 & Tests et automatisation & \code{make test} vert\\
S5 & Correction et recette & 3/3 tests avec preuves\\
S6 & Documentation et soutenance & Dossier final livré\\
\end{tabular}
\end{center}

\end{document}
